% 分类目录：ml
\documentclass[12pt, a4paper]{article}
\usepackage[margin=2.5cm]{geometry}
\usepackage{ctex}
\usepackage{amsmath, amssymb}
\usepackage{listings}
\usepackage{xcolor}
\usepackage{tabularx}
\usepackage{hyperref}

\lstset{
    language=Python,
    basicstyle=\ttfamily\small,
    keywordstyle=\color{blue},
    commentstyle=\color{green!50!black},
    stringstyle=\color{red},
    showstringspaces=false,
    breaklines=true,
    numbers=left,
    numberstyle=\tiny\color{gray},
    frame=single
}

\title{知识点：线性代数基础 (Linear Algebra Foundations)}
\author{}
\date{}

\begin{document}
\maketitle

\section{核心对象}
线性代数主要研究以下四类数学对象：
\begin{itemize}
    \item \textbf{标量 (Scalar)}：单个数字（如 $a \in \mathbb{R}$）。
    \item \textbf{向量 (Vector)}：一列有序数字（如 $\mathbf{v} \in \mathbb{R}^n$）。机器学习通常将其视为空间中的坐标点。
    \item \textbf{矩阵 (Matrix)}：二维数组（如 $\mathbf{A} \in \mathbb{R}^{m \times n}$）。表示线性映射（Linear Map）。
    \item \textbf{张量 (Tensor)}：多维数组。标量、向量、矩阵分别可看作 0、1、2 阶张量。
\end{itemize}

\section{核心运算}

\subsection{矩阵乘法}
对于 $\mathbf{C} = \mathbf{AB}$，其中 $\mathbf{A} \in \mathbb{R}^{m \times k}, \mathbf{B} \in \mathbb{R}^{k \times n}$，则 $\mathbf{C} \in \mathbb{R}^{m \times n}$：
$$ C_{ij} = \sum_{l=1}^k A_{il} B_{lj} $$
\textbf{注意}：矩阵乘法不满足交换律（$\mathbf{AB} \ne \mathbf{BA}$），但满足结合律。

\subsection{逆矩阵与行列式}
\begin{itemize}
    \item \textbf{逆矩阵} $\mathbf{A}^{-1}$：满足 $\mathbf{A} \mathbf{A}^{-1} = \mathbf{I}$。仅当 \textbf{行列式} $\det(\mathbf{A}) \ne 0$ 时存在。
    \item \textbf{行列式} $\det(\mathbf{A})$：衡量矩阵所代表的线性变换对空间体积的缩放比例。
\end{itemize}

\section{范数 (Norms)}
用于衡量向量或矩阵的“大小”，在正则化中至关重要：
\begin{enumerate}
    \item $L^1$ \textbf{范数}：$\|\mathbf{x}\|_1 = \sum |x_i|$。对应 Lasso 正则化，产生稀疏性。
    \item $L^2$ \textbf{范数}：$\|\mathbf{x}\|_2 = \sqrt{\sum x_i^2}$。最常用，对应欧氏距离。
    \item \textbf{Frobenius 范数}：衡量矩阵大小，$\|\mathbf{A}\|_F = \sqrt{\sum_{i,j} A_{ij}^2}$。
\end{enumerate}

\section{特征分解 (Eigen-decomposition)}
对于方阵 $\mathbf{A}$，若满足 $\mathbf{A}\mathbf{v} = \lambda \mathbf{v}$，则称 $\mathbf{v}$ 为 \textbf{特征向量}，$\lambda$ 为 \textbf{特征值}。
\\ \textbf{直观理解}：矩阵 $\mathbf{A}$ 对向量 $\mathbf{v}$ 的变换表现为简单的缩放。这一性质是主成分分析 (PCA) 的数学基础。

\section{奇异值分解 (SVD)}
这是线性代数最伟大的工具之一，适用于任何 $m \times n$ 的矩阵 $\mathbf{A}$：
$$ \mathbf{A} = \mathbf{U}\mathbf{\Sigma}\mathbf{V}^T $$
其中 $\mathbf{Sigma}$ 是对角阵（对角线元素为奇异值）。
\\ \textbf{应用}：数据压缩、推荐系统、降低数据噪声。

\section{特殊矩阵}
\begin{itemize}
    \item \textbf{单位矩阵} ($\mathbf{I}$)：对角线为 1，其余为 0。
    \item \textbf{对称矩阵}：$\mathbf{A} = \mathbf{A}^T$。特征值均为实数且特征向量相互正交。
    \item \textbf{正交矩阵}：$\mathbf{Q}^T\mathbf{Q} = \mathbf{I}$，即 $\mathbf{Q}^{-1} = \mathbf{Q}^T$。表示纯旋转/翻转变换，不改变长度。
\end{itemize}

\section{线性相关性与秩 (Rank)}
\begin{itemize}
    \item \textbf{线性无关}：一组向量中没有任何一只可以被其他向量线性表示。
    \item \textbf{秩 (Rank)}：矩阵中线性无关的行数或列数，代表矩阵所包含信息的维度。
\end{itemize}

\section{关键代码片段 (NumPy)}
\begin{lstlisting}
import numpy as np

A = np.array([[1, 2], [3, 4]])
# 1. 矩阵乘法
C = A @ A.T

# 2. 特征分解
eigenvalues, eigenvectors = np.linalg.eig(A)

# 3. SVD 分解
U, S, Vh = np.linalg.svd(A)
\end{lstlisting}

\section{深度面试题}

\textbf{问：为什么在机器学习中经常使用 $L^2$ 范数的平方而非 $L^2$ 范数本身？}\\
答：因为 $L^2$ 范数包含开方运算，其在原点附近的导数并不平滑。而其平方 $\|\mathbf{x}\|_2^2 = \sum x_i^2$ 处处可微且导数形式极简（$2x_i$），在利用梯度下降更新参数时计算更高效、鲁棒。

\textbf{问：正交矩阵 (Orthogonal Matrix) 有什么独特的优势？}\\
答：正交矩阵的逆等于其转置，这使得求逆运算极其简单。此外，正交变换保持向量的内积（角度）和范数（长度）不变，常用于保证数值计算的稳定性（如 QR 分解）。

\section{参考文献}
\begin{enumerate}
    \item Strang, G. (2016). \textit{Introduction to Linear Algebra}. Wellesley-Cambridge Press.
    \item Goodfellow, I., Bengio, Y., \& Courville, A. (2016). \textit{Deep Learning}. MIT Press (Chapter 2: Linear Algebra).
    \item Boyd, S., \& Vandenberghe, L. (2018). \textit{Introduction to Applied Linear Algebra}. Cambridge University Press.
\end{enumerate}

\end{document}
